\documentclass[11pt]{article}
\usepackage{fontspec,amsmath,amssymb,amsthm,hyperref}
\usepackage[margin=1in]{geometry}
\setlength{\emergencystretch}{3em}
\title{Derive the chi-square positive-part tail from the Gaussian density}
\author{LangevinLean mathematical exposition}
\date{}
\begin{document}\maketitle
\noindent Source module: \texttt{KineticLangevin/GaussianTail.lean}. Compile this proof document with XeLaTeX or LuaLaTeX. The accompanying HTML pairs the mathematical argument with the exact Lean statements and records the scope of the formalization.


The Hessian comparison needs a Gaussian tail estimate even when the Gaussian velocity and position discrepancy are dependent. The relevant bound is \(\mathbb E(\lVert G\rVert^2-a)_+\le4e^{-(a-p)/4}\) for \(a\ge5p\). The proof below derives its exponential moment and the tail estimate from actual Gaussian measures.


\section{Compute the scalar exponential square moment with its integrability}


Multiply the standard Gaussian density by \(e^{x^2/3}\). Its exponent becomes \(-x^2/2+x^2/3=-x^2/6\). The real Gaussian integral therefore gives \(\mathbb Ee^{G^2/3}=\sqrt3\). The proof establishes integrability of this weighted density before computing its integral, so the Bochner integral cannot conceal divergence. Lean uses the actual Gaussian density, the with-density integrability equivalence, the real Gaussian integral and exact square-root algebra.



\noindent Lean declarations: \texttt{scalarGaussian\_exp\_square\_third}.

\section{Multiply the independent coordinate moments}


Since \(e^{\lVert G\rVert^2/3}=\prod_i e^{G_i^2/3}\), product integrability and finite-product Fubini give \(\mathbb Ee^{\lVert G\rVert^2/3}=(\sqrt3)^p\). The proof works on the constructed product of scalar Gaussian measures and uses its established Euclidean standard-Gaussian pushforward. This handles dimension zero as well.



\noindent Lean declarations: \texttt{stdGaussian\_exp\_norm\_square\_third}.

\section{Choose an elementary Chernoff tilt}


The elementary bound \(e^u\ge1+u\) at \(u=1/3\), raised to the fourth power, proves \(\sqrt3\le e^{2/3}\). Also \((y-a)_+\le3e^{(y-a)/3}\), since its positive branch is bounded using the same exponential inequality and its zero branch is nonnegative. These two inequalities give a convenient exponential moment estimate without importing an unproved chi-square concentration result.



\noindent Lean declarations: \texttt{sqrt\_three\_le\_exp\_two\_thirds}, \texttt{positivePart\_le\_exp\_third}.

\section{Integrate the positive-part bound and use the threshold}


Applying the pointwise bound to \(y=\lVert G\rVert^2\) gives \(\mathbb E(y-a)_+\le3e^{-a/3}(\sqrt3)^p\le3e^{-a/3+2p/3}\). If \(a\ge5p\), the exponent is at most \(-(a-p)/4\). Replacing the factor three by four yields exactly the source bound. Lean first derives integrability of the positive part from its integrable exponential majorant, then applies integral monotonicity, the exact moment and scalar exponent comparison.



\noindent Lean declarations: \texttt{stdGaussian\_norm\_square\_positivePart\_le}.

\section{Use the bound on the original coupled probability space}


The pushforward integration identity transfers the estimate to any actual standard-Gaussian velocity. Its joint distribution with a numerical position discrepancy is unrestricted. This is essential for a synchronous sampling proof, where past noise generally makes those two variables dependent.



\noindent Lean declarations: \texttt{gaussian\_norm\_square\_positivePart\_le}.

\end{document}
